\emph{Proof of the local ACE benchmark.}
Fix \(n\).  The result \(\text{lem:jms-ace-theorem-five-four}\) is stated for every positive value of its cumulant-separation parameter.  Apply it with the numerical substitution \(\delta=\delta_n\).  Under this substitution the model class in \(\text{def:jms-ace-class}\) is exactly \(\mathcal C_n(\delta_n)\), because its only altered restriction is
\[
 \lvert\kappa_{r+1}(\eta)\rvert\geq\delta
 \quad\longmapsto\quad
 \lvert\kappa_{r+1}(\eta)\rvert\geq\delta_n.
\]
The quantities \(a_{1,n}\), \(b_{1,n}\), and \(a_2\) do not involve \(\delta\), whereas \(b_{2,n}\) becomes
\[
 \widetilde b_{2,n}
 =200\min\{1,C_{\theta}\}\delta_n
 \max\!\left\{
   \varepsilon_{1,n},\varepsilon_{2,n},
   (\gamma n)^{-1/2}(\psi_{\xi}+C_{\theta}\psi_{\eta})
 \right\}^{-1}.
\]
Consequently, the displayed local eligibility condition in the theorem is precisely the eligibility condition of \(\text{lem:jms-ace-theorem-five-four}\) after this substitution.  The hypotheses \(\varepsilon_{1,n},\varepsilon_{2,n}>0\), together with the requirement that all displayed logarithms and denominators be defined, supply the positivity and boundary conditions required by that result.  It follows that
\[
\begin{aligned}
 &\sup_{P\in\mathcal C_n(\delta_n)}
 Q_{P,1-\gamma}
 \!\left(\left|\widehat\theta_{\mathrm{ACE},r,n}-\theta_0(P)\right|\right)\\
 &\quad\leq
 C_{\gamma}r!16^r\delta_n^{-1}
 \Bigl[
   \varepsilon_{1,n}^r\varepsilon_{2,n}
   +C_{\theta}\varepsilon_{1,n}^{r+1}\\
 &\hspace{42mm}
   +64(C_g+\psi_{\eta})^r
    \{r^2(C_g+\psi_{\eta})+\psi_{\xi}+C_{\theta}\psi_{\eta}\}
    (\gamma n)^{-1/2}
 \Bigr].
\end{aligned}
\]
This is the claimed quantity \(U_{\mathrm{ACE},n}^{\mathrm{loc}}\).  A separately valid DML bound and this ACE bound concern two procedures under their respective hypotheses.  Their pointwise minimum is therefore an oracle upper envelope only: the two inequalities alone neither define a data-driven selector nor furnish a lower bound for \(\mathcal C_n(\delta_n)\).  Thus they imply no minimax phase diagram.

\emph{Proof for the Gaussian--Rademacher path.}
Now set \(k=4\), so \(r=3\), and fix \(0<a\leq1\).  Independence of \(G_{\circ}\) and \(S\) gives, for every \(z\in\mathbb C\),
\[
 M_a(z)
 =E[e^{z\sqrt{1-a^2}G_{\circ}}]E[e^{zaS}]
 =\exp\!\left(\frac{1-a^2}{2}z^2\right)\cosh(az).
\]
Cumulants add for independent sums and are homogeneous under scaling.  The fourth Gaussian cumulant is zero, while
\[
 \kappa_4(S)=E[S^4]-3E[S^2]^2=1-3=-2.
\]
Hence \(\kappa_4(\eta_a)=-2a^4\).  The characteristic function is
\[
 \varphi_a(t)=M_a(it)
 =\exp\!\left(-\frac{1-a^2}{2}t^2\right)\cos(at).
\]
Its exponential factor is strictly positive for real \(t\), and therefore its first positive zero is \(t_a=\pi/(2a)\).  For every \(d\in\mathbb R\),
\[
 E\!\left[\sin\{t_a(\eta_a+d)\}\right]
 =\operatorname{Im}\{e^{it_ad}\varphi_a(t_a)\}=0.
\]

We next prove the relevance and positivity statement.  Direct differentiation at the zero yields
\[
 \varphi_a'(t_a)
 =-a\exp\!\left\{-\frac{(1-a^2)\pi^2}{8a^2}\right\}
 =-A_a.
\]
The variable \(\eta_a\) has moments of every order, so differentiation under the expectation is justified.  For every fixed \(d\in\mathbb R\),
\[
\begin{aligned}
 E[(\eta_a+d)\sin\{t_a(\eta_a+d)\}]
 &=-\operatorname{Re}\!\left[
   \left.\frac{d}{dt}\{e^{itd}\varphi_a(t)\}\right|_{t=t_a}
 \right]\\
 &=-\operatorname{Re}\{e^{it_ad}\varphi_a'(t_a)\}
 =A_a\cos(t_ad).
\end{aligned}
\]
Because \(\eta_a\perp X\) and \(D_n\) is measurable with respect to \(X\), conditioning on \(D_n\) and integrating gives
\[
 m_a:=E_P[Z_n\sin(t_aZ_n)]
 =A_aE_P[\cos(t_aD_n)].
\]
Here \(s\geq r+1=4\).  Lyapunov's inequality and \(\text{ass:treatment-code-radius}\) therefore imply
\[
 E_P|D_n|
 \leq \|D_n\|_{L^s(P_X)}
 \leq \varepsilon_{1,n}.
\]
Using \(1-\cos u\leq |u|\) and the assumed condition \(t_a\varepsilon_{1,n}\leq1/2\), we obtain
\[
 E_P[\cos(t_aD_n)]
 \geq1-t_aE_P|D_n|
 \geq1-t_a\varepsilon_{1,n}
 \geq\frac12.
\]
Thus \(m_a\geq A_a/2>0\), proving both the identity and the denominator positivity required by the ratio.

It remains to establish the uniform risk bound.  Define
\[
 W=Z_n\sin(t_aZ_n),
 \qquad
 V=Y\sin(t_aZ_n),
 \qquad
 R=(Y-\theta_0Z_n)\sin(t_aZ_n).
\]
The residual representation is \(Y-\theta_0Z_n=b_n(X)+\xi\).  Conditional shift annihilation and \(\eta_a\perp X\) give
\[
 E_P[b_n(X)\sin(t_aZ_n)]
 =E_P\!\left[b_n(X)E_P\{\sin(t_aZ_n)\mid X\}\right]=0.
\]
Also \(\sin(t_aZ_n)\) is measurable with respect to \((X,T)\), so \(\text{ass:outcome-mean-independence}\) gives
\[
 E_P[\xi\sin(t_aZ_n)]
 =E_P\!\left[\sin(t_aZ_n)E_P(\xi\mid X,T)\right]=0.
\]
Consequently \(E_P[R]=0\) and \(E_P[V]=\theta_0m_a\).  This is the identifying population moment; it is distinct from the sample estimator in the theorem.

We now verify uniform square integrability.  Clipping and \(\text{ass:g-range}\) imply \(|D_n|\leq2C_g\).  Since \(E[\eta_a^2]=1\), \(E[\eta_a]=0\), and \(\eta_a\perp D_n\),
\[
 E_P[W^2]
 \leq E_P[Z_n^2]
 \leq 1+4C_g^2
 =:K_W.
\]
Moreover, \(|b_n(X)|\leq C_q+2C_{\theta}C_g\).  The conditional sub-Gaussian restriction implies \(E_P[\xi^2]\leq c_{\psi}\psi_{\xi}^2\), where \(c_{\psi}\) is the universal constant fixed by the convention for the \(\psi_2\)-norm.  Therefore
\[
 E_P[R^2]
 \leq 2(C_q+2C_{\theta}C_g)^2+2c_{\psi}\psi_{\xi}^2
 =:K_R.
\]
Both constants are finite and independent of \(a\), \(n\), and the particular path law.

Write \(\overline W=\mathbb P_nW\), \(\overline V=\mathbb P_nV\), and \(\mathcal E=\{\overline W\geq A_a/4\}\).  On \(\mathcal E\), projection onto \([-C_{\theta},C_{\theta}]\) is a contraction and \(\theta_0\) belongs to this interval.  Hence
\[
 \left|\widehat\theta_{a,n}-\theta_0\right|
 \leq\left|\frac{\overline V}{\overline W}-\theta_0\right|
 =\frac{|\mathbb P_nR|}{\overline W}
 \leq\frac{4}{A_a}|\mathbb P_nR|.
\]
By \(\text{ass:iid-sampling}\) and \(E_P[R]=0\),
\[
 E_P[(\mathbb P_nR)^2]
 =\frac{\operatorname{Var}_P(R)}{n}
 \leq\frac{K_R}{n}.
\]
On \(\mathcal E^c\), the estimator equals zero, so its squared error is at most \(C_{\theta}^2\).  Since \(m_a\geq A_a/2\), Chebyshev's inequality and the same sampling assumption yield
\[
 P(\mathcal E^c)
 \leq P\!\left(|\overline W-m_a|\geq\frac{A_a}{4}\right)
 \leq\frac{16K_W}{nA_a^2}.
\]
Combining these bounds gives
\[
 E_P[(\widehat\theta_{a,n}-\theta_0)^2]
 \leq\frac{16(K_R+C_{\theta}^2K_W)}{nA_a^2}.
\]
Independently, both the estimator and the target belong to \([-C_{\theta},C_{\theta}]\), so the same risk is at most \(4C_{\theta}^2\).  Taking the smaller bound and using
\[
 A_a^{-2}
 =a^{-2}\exp\!\left\{\frac{(1-a^2)\pi^2}{4a^2}\right\}
\]
proves the asserted inequality with a finite constant depending only on \(C_{\theta},C_g,C_q,\psi_{\xi}\).

Finally, \(\Delta_a=2a^4\) implies
\[
 a=2^{-1/4}\Delta_a^{1/4},
 \qquad
 a^{-2}=\sqrt2\,\Delta_a^{-1/2}.
\]
The condition \(t_a\varepsilon_{1,n}\leq1/2\) is therefore equivalent to
\[
 \varepsilon_{1,n}
 \leq\frac{a}{\pi}
 =\frac{\Delta_a^{1/4}}{2^{1/4}\pi}.
\]
Furthermore,
\[
 a^{-2}\exp\!\left\{\frac{(1-a^2)\pi^2}{4a^2}\right\}
 =\sqrt2e^{-\pi^2/4}\Delta_a^{-1/2}
 \exp\!\left\{\frac{\pi^2\sqrt2}{4}\Delta_a^{-1/2}\right\}.
\]
Absorbing the fixed factor \(\sqrt2e^{-\pi^2/4}\) into \(C\) proves the final risk display.  This argument uses only the explicit Gaussian--Rademacher path and supplies no converse over the larger class defined solely by fourth-cumulant separation.  It therefore proves the stated construction-specific benchmark, not a full minimax phase diagram.